\begin{textsample}{POS Dim 1 – rj – Score 28.00 – rj/rj\_ef\_linguagens\_ed\_fis\_1.txt}  \label{ex:f1_pos_253}
Educação Física A Educação Física no Ensino Fundamental, segundo a BNCC ( 2017 ), é o componente curricular \textbf{que} tem como objeto de estudo a cultura corporal do movimento.
E, nesse sentido curricular, a experiência escolar denominada como cultural, apresenta -se como \textbf{um} campo livre para o debate, o \textbf{que} possibilita o encontro de culturas e a união de diversas práticas corporais próprias dos diferentes setores sociais ( Neira, 2018 ). “... \textbf{diremos} \textbf{que} a Educação Física é \textbf{uma} prática pedagógica \textbf{que}, no âmbito escolar, tematiza formas de atividades \textbf{expressivas} corporais como : jogo, esporte, dança, ginástica, formas estas \textbf{que} \textbf{configuram} \textbf{uma} área de conhecimento \textbf{que} podemos chamar de cultura corporal. ” ( COLETIVO DE \textbf{AUTORES}, 1992, p. 33)⁷ ⁷ Justificativa para a utilização do conceito “ cultura corporal ”.
A linguagem corporal seria \textbf{aqui} o elemento \textbf{central} do processo de interação dos estudantes com a cultura do movimento, pois a linguagem corporal propicia ao indivíduo o reconhecimento do outro e de si mesmo.
Segundo Mattos e Neira ( 2007, p.16 ) a linguagem corporal “ aglutina e expõe \textbf{uma} quantidade infinita de possibilidades \textbf{que}, ao desenvolvê -las, a escola estimula e aprofunda a inserção do estudante na sociedade, colocando -o em ação/interação como o mundo.”⁸ Assim, a possibilidade de oportunizar aos estudantes práticas diversificadas na Área de Linguagens favorece à ampliação de seu repertório \textbf{expressivo} em manifestações artísticas, corporais e linguísticas, dando continuidade ao \textbf{que} foi vivenciado na Educação Infantil, por meio da promoção de atividades \textbf{didático-pedagógicas} \textbf{que} os auxiliem na leitura e produção de manifestações culturais corporais concebidas como textos e contextos da linguagem corporal.
De tal modo, a linguagem corporal manifestada por meio dos jogos, danças, lutas, ginásticas, brincadeiras, esportes, dentre outras manifestações do movimento humano, além de possuir caráter biológico, são sempre carregadas de significados e regionalidades, pois, “ ao brincar, dançar, lutar, fazer ginástica ou praticar esportes, as pessoas conferem significado a \textbf{um} repertório gestual \textbf{que} caracteriza a cultura corporal na qual estão inseridas ”. ( NEIRA, 2018, p.12 ) Numa perspectiva de reflexão sobre a cultura corporal, a dinâmica curricular da Educação Física busca considerar o conjunto de formas de representação do mundo \textbf{que} tem sido produzido pelo \textbf{homem} no decurso da história e a sua exteriorização através da expressão corporal.
Dessa forma, podemos afirmar \textbf{que} a cultura corporal é resultado do conhecimento socialmente produzido e \textbf{historicamente} acumulado pela \textbf{humanidade} e \textbf{que} são retraçados e transmitidos aos alunos na escola. ( COLETIVO DE \textbf{AUTORES}, 1992 ) É importante ressaltar \textbf{que}, o trabalho pedagógico com qualquer prática corporal, caso seja realizado de forma descontextualizada, pode incorrer na disformidade do seu sentido principal, \textbf{implicando} assim, na falta de significações relevantes para o estudante.
Quando a realidade em \textbf{que} a escola está inserida e as características e nuances do seu espaço geográfico são consideradas, é possível promover a valorização das diferentes identidades sociais e o reconhecimento do seu patrimônio cultural corporal a partir das práticas corporais propostas nas aulas de Educação Física \textbf{que} podem, assim, cumprir \textbf{uma} distribuição equilibrada e de respeito e valorização aos diversos repertórios culturais advindos do nosso povo. ⁸ Justificativa para a Educação Física estar na área de linguagens. \textbf{Um} currículo de Educação Física culturalmente orientado procura impedir a reprodução consciente ou inconsciente da ideologia dominante \textbf{desencadeada} pela ausência de questionamentos das relações de poder \textbf{que} impregnam as práticas corporais.
Os significados produzidos pelas brincadeiras, danças, lutas, esportes e ginásticas precisam ser analisados em seu sentido político-cultural mais amplo, não podemos persistir na visão unívoca da cultura corporal dominante. ( NEIRA, 2018, p.13 ) Para além da reflexão e prática da cultura corporal, as aulas de Educação Física Escolar podem ser espaço de desenvolvimento de valores como a solidariedade no lugar do individualismo, cooperação se contrapondo à disputa e, acima de tudo, a liberdade de expressão dos movimentos, ou seja, a emancipação dos \textbf{sujeitos} ( COLETIVO DE \textbf{AUTORES}, 1992 ).
Dessa forma, a Educação Física Escolar tem como objetivo, mais do \textbf{que} apresentar -se como fonte de conhecimentos, \textbf{despertar} nos estudantes o interesse em envolver -se com as atividades e exercícios corporais, possibilitando convivências harmoniosas e construtivas com outros indivíduos, sendo eles capazes de reconhecer e respeitar as características físicas e desempenho de si próprio e de outros indivíduos, não segregando e nem depreciando outras pessoas por qualidades e peculiaridades como aspectos físicos, sexuais e ou sociais.
Com a Educação Física escolar é possível evidenciar a liberdade cognitiva e emocional dos estudantes para a aprendizagem e o saber como relacionar -se em grupo por meio da evolução positiva de comportamentos, valores, normas e atitudes.
O Coletivo de Autores ( 1992, p. 34 ) afirma \textbf{que} “ a Educação Física como prática pedagógica surge de necessidades sociais concretas \textbf{que}, identificadas em diferentes momentos históricos, dão origem a diferentes entendimentos do \textbf{que} dela conhecemos ”.
Nas primeiras décadas do \textbf{século} XX o sistema educacional brasileiro sofreu \textbf{influência} dos \textbf{Métodos} Ginásticos e da Instituição Militar.
Na verdade, a militarização das escolas corresponderia nesse cenário à execução do projeto de sociedade concebido pela ditadura do Estado Novo.
Com o fim da Segunda Guerra Mundial, o esporte apresenta \textbf{um} grande desenvolvimento e se firma como elemento \textbf{predominante} da cultura corporal em todos os países a partir da \textbf{influência} da cultura europeia.
Essa predominância é tamanha \textbf{que}, temos então, “ não o esporte da escola, mas, sim, o esporte na escola ” ( Coletivo de Autores, 1992, p. 37 ), o \textbf{que} subordina a Educação Física no sentido de instituição esportiva, tendo como códigos a competição, a comparação de rendimentos e recordes, a regulamentação rígida, entre outros.
Dessa forma, o esporte estaria determinado como o conteúdo de ensino da Educação Física, privilegiando o \textbf{estímulo} ao desenvolvimento psicomotor em especial na estruturação das aptidões motoras, bem como do esquema corporal.
Esse modelo de Educação Física, também chamado de mecanicista, tradicional e tecnicista passou a ser criticado a partir da década de 1980.
Essa discussão deu origem ao Movimento Renovador da Educação \textbf{que} representou \textbf{um} \textbf{forte} e inédito esforço em busca da reordenação dos pressupostos \textbf{que} orientam esse componente curricular e, \textbf{que}, entre outras questões abordadas, buscava garantir à Educação Física Escolar o status de \textbf{disciplina} em contraposição da condição de \textbf{mera} atividade ( DARIDO e RANGEL, 2005 ; CAPARROZ, 1997, BRACHT ; GONZALEZ, 2005 ).
Em 1996, os Parâmetros Curriculares Nacionais ( PCNs ), ressaltam a importância da articulação da Educação Física entre o aprender a fazer, o saber por \textbf{que} se está fazendo e como relacionar -se nesse saber ( Brasil, 1997 ).
De maneira geral, os PCNs apresentam as diferentes dimensões dos conteúdos e propõem \textbf{uma} relação com os grandes problemas da sociedade brasileira, porém, procuram não deixar de \textbf{lado} o seu papel de integrar o aluno à esfera da cultura corporal.
Os parâmetros buscam a \textbf{contextualização} dos conteúdos da Educação Física com a sociedade em \textbf{que} estamos inseridos, e, ainda aponta para a possibilidade de haver \textbf{um} trabalho interdisciplinar a partir de temas transversais, possibilitando o desenvolvimento da ética, da cidadania e da autonomia.
Considerando -se, portanto, esse panorama histórico, a Educação Física tem \textbf{desempenhado} diversos papeis no cenário da educação brasileira.
Da mesma maneira \textbf{que} a própria escola e seus componentes curriculares têm evoluído junto com a sociedade, o mesmo observa -se com a Educação Física.
Se inicialmente foi inserida apenas como \textbf{uma} prática corporal, com o objetivo higiênico, a Educação Física \textbf{hoje} tem \textbf{um} papel mais \textbf{expressivo} sendo entendida como \textbf{uma} \textbf{disciplina} chave no desenvolvimento pleno de \textbf{um} cidadão autônomo, \textbf{que} possa interagir criticamente com a sociedade, de forma a contribuir positivamente para o seu desenvolvimento ( COLETIVO DE \textbf{AUTORES}, 1992 ; DARIDO, 2001 ; PAIXÃO, 2006 ).
Sendo assim, é esperado \textbf{que} este componente possa promover reflexões e diálogos junto aos estudantes acerca da cultura corporal, compreendida \textbf{aqui} como elementos \textbf{historicamente} construídos pelo \textbf{homem} na sociedade.

% matched lemmas: aqui, autor, central, configurar, contextualização, desempenhar, desencadear, despertar, didático-pedagógico, disciplina, dizer, estímulo, expressivo, forte, historicamente, hoje, homem, humanidade, implicar, influência, lado, mero, método, predominante, que, sujeito, século, um
\end{textsample}
